\section{Morfismos cuasifinitos}
\label{I.2}
\begin{proposition}
\label{I.2.1}
Sea $A\to B$ un homomorfismo local (N.B.: los anillos son ahora noetherianos), y sea \ifthenelse{\boolean{orig}}{$m$}{$\goth{m}$} el ideal maximal de~$A$. Son equivalentes las condiciones siguientes:
\begin{enumerate}
\item[(i)] \ifthenelse{\boolean{orig}}
{$B/mB$ es de dimensión finita sobre~$k=A/m$}
{$B/\goth{m}B$ es de dimensión finita sobre~$k=A/\goth{m}$}.
\item[(ii)] \ifthenelse{\boolean{orig}}
{$mB$ es un ideal de definición, y $B/r(B)=\kappa(B)$ es una extensión de~$k=\kappa(A)$}
{$\goth{m}B$ es un ideal de definición, y $B/\goth{r}(B)=\kres(B)$ es una extensión de~$k=\kres(A)$}.
\item[(iii)] El completado $\widehat{B}$ es finito sobre el completado $\widehat{A}$ de~$A$.
\end{enumerate}
\end{proposition}
Se dice entonces
\marginpar{2}
que $B$ \emph{es cuasifinito}
\index{cuasifinito (morfismo, álgebra)|hyperpage}%
sobre~$A$. Se dice que un morfismo $f\colon X\to Y$ es cuasifinito en~$x$ (o que el $Y$-preesquema $f$ es cuasifinito en~$x$) si $\cal{O}_x$ es cuasifinito sobre~$\cal{O}_{f(x)}$. Esto equivale asimismo a decir que $x$ está \emph{aislado en su fibra}~$f^{-1}(f(x))$.% SGA1-CORR-0003: diplomatic French and original print have f^{-1}(x), with x in X.
Se dice que un morfismo es cuasifinito si lo es en todo punto\footnote{En EGA~II~6.2.3 se supone además que $f$ es de tipo finito.}.
\begin{corollaire}
\label{I.2.2}
Si $A$ es completo, ser cuasifinito equivale a ser finito.
\ifthenelse{\boolean{orig}}
{Podría darse el sorites habitual \textup{(i) (ii) (iii) (iv) (v)} para los morfismos cuasifinitos, pero no parece indispensable aquí.}
{}
\end{corollaire}

\ifthenelse{\boolean{orig}}
{}
{Podría darse el sorites habitual \textup{(i), (ii), (iii), (iv), (v),} para los morfismos cuasifinitos, pero no parece indispensable aquí.}

